\documentclass{../class/twocolpaper}

\usepackage{amsmath,amssymb}
\usepackage{graphicx}
\usepackage{booktabs}
\usepackage[numbers,sort&compress]{natbib}
\usepackage{hyperref}

\title{A Simple Example of a Two-Column Academic Paper}

\author{
    Author A$^{1,*}$, Author B$^{1,2}$, Author C$^{2}$
    \\
    \small
    \parbox{0.95\linewidth}{
        \centering
        $^{1}$ \itshape Affiliation, University, Address, City, Postcode, State/Province, Country
        \\
        $^{2}$ \itshape Affiliation, University, Address, City, Postcode, State/Province, Country
        \\
        $^{*}$ Corresponding author: author@example.com
    }
}

\date{}

\begin{document}

\maketitle

\begin{abstract}
This paper provides a simple example of using the
\texttt{twocolpaper} document class. The template is designed
for academic manuscripts in a compact two-column format.
The abstract spans both columns, while the main text is
automatically arranged in two columns.
\end{abstract}

\section{Introduction}

Two-column layouts are widely used in academic publications
because they provide a compact presentation of scientific
content. This example demonstrates the basic usage of the
\texttt{twocolpaper} class.

The document class is based on the standard \texttt{article}
class and provides a fixed A4 two-column layout. Users can
write a manuscript using the standard \LaTeX{} commands without
having to manually configure the page geometry.

\section{Equation}

Consider a simple quadratic function,

\begin{equation}
    y = ax^2 + bx + c.
    \label{eq:quadratic}
\end{equation}

For $a=1$, $b=2$, and $c=1$, Eq.~\ref{eq:quadratic} becomes
$y=x^2+2x+1$.

\section{Table}

Table~\ref{tab:example} gives several example values calculated
from Eq.~\ref{eq:quadratic}.

\begin{table}[htbp]
    \centering
    \caption{Example values calculated from Eq.~\ref{eq:quadratic}.}
    \label{tab:example}
    \begin{tabular}{ccccc}
        \toprule
        $a$ & $b$ & $c$ & $x$ & $y$ \\
        \midrule
        1.0 & 2.0 & 1.0 & $-2$ & 1 \\
        1.0 & 2.0 & 1.0 & $-1$ & 0 \\
        1.0 & 2.0 & 1.0 & 0   & 1 \\
        1.0 & 2.0 & 1.0 & 1   & 4 \\
        1.0 & 2.0 & 1.0 & 2   & 9 \\
        \bottomrule
    \end{tabular}
\end{table}

\section{Figure}

Figure~\ref{fig:example} illustrates an example figure
placeholder.

\begin{figure}[htbp]
    \centering
    \fbox{\rule{0pt}{40mm}\rule{0.85\columnwidth}{0pt}}
    \caption{Example figure.}
    \label{fig:example}
\end{figure}


\section{Citation}

Mathematical typesetting and document preparation are important
components of scientific writing~\cite{lamport1994}.

\bibliographystyle{plain}
\bibliography{references}

\end{document}